\documentclass{article}

\usepackage[utf8]{inputenc}
\usepackage{helvet}
\usepackage{geometry}
\usepackage{xcolor}
\usepackage{eso-pic}
\usepackage{fancyhdr}
\usepackage{parskip}
\usepackage{microtype}
\usepackage[english, ukrainian]{babel}
\usepackage{url}
\usepackage{amsmath}
\usepackage{amssymb}
\usepackage{amsthm}
\usepackage{longtable}
\usepackage{booktabs}
\usepackage{hyperref}
\usepackage{titlesec}

\definecolor{nato}{HTML}{293757}
\definecolor{footergray}{gray}{0.65}

\geometry{a4paper,left=3cm,right=3cm,top=3cm,bottom=3cm}
\renewcommand{\familydefault}{\sfdefault}
\setcounter{secnumdepth}{3}

\DeclareFontFamily{T2A}{phv}{}
\DeclareFontShape{T2A}{phv}{m}{n}{<->ssub * cmss/m/n}{}
\DeclareFontShape{T2A}{phv}{bx}{n}{<->ssub * cmss/bx/n}{}
\DeclareFontShape{T2A}{phv}{m}{it}{<->ssub * cmss/m/it}{}
\DeclareFontShape{T2A}{phv}{bx}{it}{<->ssub * cmss/bx/it}{}

\titleformat{\section} {\color{nato}\Huge\bfseries}{\thesection.}{1em}{}
\titlespacing*{\section}{0pt}{30pt}{12pt}
\titleformat{\subsection} {\color{nato}\LARGE\bfseries}{\thesubsection.}{1em}{}
\titlespacing*{\subsection}{0pt}{20pt}{8pt}
\titleformat{\subsubsection} {\color{nato}\Large\bfseries}{\thesubsubsection.}{1em}{}
\titlespacing*{\subsubsection}{0pt}{10pt}{6pt}

\fancyhf{}
\fancyhead[R]{\small\color{footergray} 27 червня 2026}
\fancyfoot[L]{\small\color{footergray} ERP/1. Модуль Конференцзв'язок}
\fancyfoot[R]{\small\color{footergray}\thepage}

\newcommand\HUGE[1]{\fontsize{40}{40}\selectfont #1}
\renewcommand{\headrulewidth}{0pt}
\renewcommand{\footrulewidth}{0.4pt}
\renewcommand{\footrule}{\color{footergray}\hrule width \textwidth height 0.4pt}

\usepackage{amsmath}
\usepackage{amssymb}
\usepackage{booktabs}
\usepackage{hyperref}
\usepackage{tikz}
\usetikzlibrary{shapes.geometric, arrows, arrows.meta, positioning, fit, backgrounds}

\let\originalfootnotesize\footnotesize
\let\oldverbatim\verbatim
\let\oldendverbatim\endverbatim
\renewenvironment{verbatim}{\originalfootnotesize\oldverbatim}{\oldendverbatim}

\title{Deterministic Real-Time Multimedia Synchronization: Architecting Low Latency Topologies in GStreamer and WebRTC}
\author{Namdak Tonpa, Zen Crypted}
\date{\today}

\begin{document}

\begin{titlepage}
  \AddToShipoutPictureBG*{\AtPageLowerLeft{\color{nato}\rule{\paperwidth}{\paperheight}}}
  \color{white}\thispagestyle{empty}\vspace*{3cm}\centering
  {\Huge  \textbf{SYNRC: OpenWebRTC} \par} \vspace{1.5cm}
  {\HUGE  \textbf{Технічне завдання} \par} \vspace{1.5cm}
  {\LARGE \textbf{для реалізації механізму синхронізації медіа потоків
                  з низьколатентними вимогами реального часу
                  і синхронізацією кадрів через головний тактовий генератор мікшера
                  } \par} \vspace{1.5cm}
  {\large На основі Технічних Вимог \par}
  {\large Згідно з Постановою КМУ № 205 від 21.02.2025 \par}
  {\large Формалізованих за стандартом ISO/IEC/IEEE 29148:2018 \par}
  \vfill
  {\large 2026 \copyright\ Криптографічні Телесистеми \par}
\end{titlepage}

\newpage
\maketitle

\begin{abstract}
Achieving ultra-low latency and stable real-time multimedia transmission in heterogeneous
environments requires rigorous architectural precision. The primary challenge lies in the
deterministic reconciliation of conflicting asynchronous domains: non-blocking WebRTC transmission,
unpredictable I/O operations (e.g., disk recording), and strict isochronous rendering (compositing).
This article presents an engineering framework for resolving pipeline starvation, timestamp dilation,
and latency accumulation by enforcing strict thread isolation, temporal queueing, and monotonic clock
synchronization within GStreamer and WebRTC topologies.
\end{abstract}

\lefthyphenmin=1
\hyphenpenalty=100
\tolerance=6000

\tableofcontents

\newpage
\section{Introduction}

The crux of real-time multimedia stability is the maintenance of a continuous, synchronized data flow
across diverse subsystems. Naive topological designs often tightly couple sinks and processing elements
within the same synchronous state space. Consequently, when unpredictable elements—such as a \texttt{filesink}
or \texttt{hlssink2}—encounter blocking I/O, backpressure propagates upstream. If an isochronous compositor
and a WebRTC sink share this state space, the compositor blocks, stalling the overarching temporal
reference. Failure to isolate these domains dictates a departure from real-time execution, leading to
catastrophic latency accumulation.

\section{Media Sync Architecture}

\subsection{Asynchronous Domain Isolation}

\subsubsection{Thread Topology and I/O Decoupling}

In GStreamer, queues define asynchronous thread boundaries. By default, data flow and state changes
execute synchronously. If a downstream element stalls, it halts the thread of upstream elements.
To guarantee localized I/O blocking or CPU spikes do not propagate backpressure, architectures
must aggressively inject queues prior to encoders, sinks, and test sources.

This decoupling isolates components into independent OS threads, ensuring the master clock and
compositor remain untethered. Furthermore, non-blocking state changes (\texttt{async=false}, \texttt{sync=false})
must be enforced on I/O-bound sinks.

\subsubsection{Temporal Queuing and Absolute Drop Policies}

Default queue implementations impose arbitrary limits on buffer counts and bytes, which are semantically
poor metrics for real-time video where frame sizes fluctuate. To enforce an absolute zero-latency
drop policy, queues must be strictly temporal.

By setting explicit temporal boundaries (\texttt{max-size-time=300000000}, \texttt{max-size-buffers=0}, \texttt{max-size-bytes=0})
coupled with a \texttt{leaky=downstream} policy, the queue is restricted to holding precisely 300 milliseconds
of actual A/V time. If processing latency exceeds this threshold, older frames are instantaneously
jettisoned. Trading recording fidelity for live stream continuity is a non-negotiable axiom of real-time design.

For base generators (\texttt{videotestsrc}, \texttt{audiotestsrc}), which produce deterministic, constant-rate data,
micro-queues (\texttt{max-size-buffers=1} for video, \texttt{max-size-buffers=5} for audio) provide sufficient
asynchronous decoupling without semantic delay, rendering the generators immune to downstream transients.

\subsection{Synchronization and Temporal Monotonicity}

\subsubsection{Mitigating Timestamp Dilation}

Timestamp dilation—perceived as a "slow motion" effect in WebRTC playback—is a classic desynchronization anomaly.
It occurs when a processing pipeline (e.g., a compositor) is constrained by CPU saturation and fails to generate
frames at its targeted framerate. If downstream filters enforce a strict high framerate (e.g., \texttt{30/1}),
the compositor continues assigning timestamps based on the theoretical rate. Consequently, one real-time
second of processing may only yield 0.5 seconds of pipeline time, a discrepancy the WebRTC receiver faithfully obeys.

To mathematically guarantee A/V synchronization, the requested framerate must be explicitly lowered
(e.g., \texttt{framerate=15/1}) to align with the hardware's guaranteed processing capacity. This ensures
that generated frames are stamped accurately (0ms, 66ms, 133ms...), spanning exactly 1.0 real-time
second, immediately eradicating the slow-motion defect.

\subsubsection{Isochronous Interpolation for Variable Ingress}

WebRTC clients dynamically throttle their egress framerates based on network congestion. Conversely,
a compositor demands a constant, unwavering stream of frames. If a client's framerate drops,
the compositor starves, blocking the pipeline.


This impedance mismatch is resolved by injecting a \texttt{videorate} element coupled with a strict \texttt{capsfilter}
immediately post-decode. \texttt{videorate} acts as an interpolation buffer, deterministically duplicating frames
to satisfy the compositor's isochronous requirements without inducing latency.

\subsection{Hardware and Resource Constraints Optimization}

Scaling software-based encoding and high-resolution compositing requires specific architectural
adaptations on resource-constrained hardware (e.g., legacy x86, ARM platforms):

\begin{enumerate}
\item \textbf{Hardware Acceleration}: Offloading DSP tasks to dedicated ASICs (e.g., \texttt{v4l2h264enc}, QuickSync, NVENC is critical for predictable latency.
\item \textbf{Spatial Subsampling}: Reducing the compositor canvas resolution exponentially decreases memory bandwidth and CPU cycles.
\item \textbf{Thread Topology Management}: Enforcing strict CPU affinity and thread pinning mitigates context-switching overhead in the OS scheduler.
\item \textbf{Jitterbuffer Tuning}: Dynamically adjusting the \texttt{webrtcbin} internal RTP jitterbuffer latency based on measured network variance minimizes baseline delay in LAN environments while preserving stability on degraded WAN links.
\end{enumerate}

\subsection{Telemetry and Diagnostics in the RTP/RTCP Layer}

Deep visibility into the RTP/RTCP session layer is paramount for maintaining stability. Utilizing the modern Rust RTP plugin (\texttt{rsrtp}), session state, network jitter, and clock skew can be profiled.

\begin{enumerate}
\item \textbf{Jitter and Latency}: Tracing incoming packet buffering evaluates if network jitter causes late arrivals to drop or monitors how out-of-order packets are reordered. Sequence number tracking explicitly identifies lost packets, correlating network health with pipeline starvation.
\item \textbf{Clock Skew \& Synchronization}: Monitoring the skew calculation between the remote sender's clock and the local system clock (\texttt{timestamping-mode=skew}) maps RTP timestamps to local presentation time, isolating A/V sync drift.
\item \textbf{RTCP Quality Feedback}: Generating Sender Reports (SR) in \texttt{rtpsend} broadcasts packet counts and NTP timestamps. \texttt{rtprecv} handles Receiver Reports (RR), tracing NACKs (retransmission requests) or PLIs (keyframe requests) triggered by packet loss when operating under the \texttt{avpf} profile.
\end{enumerate}

\section{Conclusion}

Architecting stable, ultra-low latency WebRTC and GStreamer pipelines mandates a shift from synchronous,
tightly-coupled topologies to asynchronous, time-bounded, and heavily isolated designs. By enforcing
strict thread boundaries, temporal queue policies, and monotonic timestamping, systems can achieve
deterministic performance even under volatile network and I/O conditions.

\end{document}
